\section{Feferman--Vaught composition and Boolean powers}

\begin{annexitem}{annex:p03-fv-boundaries}{thm:p03-fv-boundaries}{Explicit separating models and the missing-choice witness family}
\begin{proof}
For the first claim take a unary relation and a product tuple whose first
coordinate satisfies it and second does not. Its local truth set is the
singleton first coordinate, not the whole two-point index set. Hence the
product relation is false and its logical negation true. The first factor
does not satisfy the negation. This proves that replacing the computed
Boolean condition for negation by universal factor negation is invalid.

In the two-element Boolean algebra, a finite partition contains exactly
one coefficient \(1\); all others are zero. Its presentation class is
therefore determined by that coefficient's one base label. This gives an
isomorphism from the base structure, since refinements compute its operations
and relations on that label. For \(B=\mathcal P(\{0,1\})\), every
two-coordinate function has finite range and allowed fibers, so the
represented-power theorem identifies it with the full two-coordinate power.
A two-element carrier has four such functions. The three functions
\((0,0),(0,1),(1,0)\) are pairwise distinct, so its three-distinct-elements
sentence is true, whereas it is false in the two-element base. Constants
give an injective structure map, but this sentence makes it non-elementary.
There is no conflict with the elementary-coefficient theorem: the two-element
Boolean algebra is not elementary in the four-element one, as the latter
satisfies \(\exists x(x\ne0\land x\ne1)\) and the former does not.

For \(A=\mathbb N\) and \(I=\mathbb N\), the identity index function
has infinite image and therefore is not in the represented finite-range
power. The preceding elementary-substructure theorem proves elementarity
of that proper subset of the full power. On the discrete topological
\(\mathbb N\), each singleton is an open constant neighborhood for the
identity function. It is locally constant, and continuous with infinite
range. The cover by all singletons has no finite subcover, proving the
missing compactness explicitly. This discrete index set is not silently
identified with the compact Stone space of \(\mathcal P(\mathbb N)\).

For the ring distinction let the base ring be
\(A=\mathbb F_2\times\mathbb F_2\) and use the two-element coefficient
algebra. Its power is isomorphic to \(A\), as just proved. This ring has
four idempotents \((0,0),(1,0),(0,1),(1,1)\), because each coordinate
is idempotent. Its idempotent Boolean algebra thus has four elements,
not the two elements of the chosen coefficient algebra. The connected-stalk
hypothesis in the source's internal ring interpretation cannot be deleted.

Finally, work conditionally in ZF with a family \((D_i)_{i\in I}\) of
nonempty sets having no choice function. Give \(M_i\) the tagged domain
\(\{(0,0)\}\cup\{(1,d):d\in D_i\}\), one constant interpreted as
\((0,0)\), and a unary predicate interpreted as the second tagged part.
The full product domain is nonempty, since the constant tuple exists.
Every factor satisfies \(\exists y\,P(y)\). But a product witness for
\(P\) would choose a \(d\in D_i\) at every index, contrary to the
assumption. The existential Boolean clause nevertheless has a valid
partition: its positive-pattern cell is \(I\), its negative cell empty,
and its local positive existential truth set is \(I\).
This is an actual failure without witness Choice, even with a nonempty
product domain. The argument does not assert that ZF proves existence of
such a family; it identifies exactly what the Choice step rules out.
\end{proof}
\end{annexitem}

\begin{annexitem}{annex:p03-fv-coordinate-terms}{thm:p03-fv-coordinate-terms}{Term induction with functions and constants}
\begin{proof}
A variable has its assigned coordinate value. A constant has its factor's
constant value by definition of the product. If the identity holds for each
immediate subterm \(t_j\), coordinatewise interpretation of \(f\) gives
\(f^M(t_1^M,\ldots,t_k^M)(i)=
f^{M_i}(t_1^{M_i},\ldots,t_k^{M_i})\), proving the compound case.
An equality of product values means equality at every coordinate; a relation
is interpreted by the same universal coordinate condition. These give exactly
the atomic truth-set criterion. For an empty index set the criterion is
vacuous, as is equality of the empty tuple values, so that case also holds.
Logical falsity is not an atomic relation and is not changed into truth by
this vacuity.
\end{proof}
\end{annexitem}
\begin{annexitem}{annex:p03-fv-worked-existential}{thm:p03-fv-worked-existential}{The two relevant witness patterns}
\begin{proof}
For the matrix, use components \(R(x,y),P(y)\) and Boolean condition
\(X_R=1\land X_P\ne1\). The existential partition has four cells,
indexed by their two truth bits. Its first condition forces the two
\(R\)-false cells to zero. Write \(Z_{10}\) for its \(R\)-true,
\(P\)-false cell and \(Z_{11}\) for its remaining cell. The requirements
become \(Z_{10}\vee Z_{11}=1\), disjointness,
\(Z_{10}\le U\), \(Z_{11}\le V\), and \(Z_{10}\ne0\).
They imply \(U\vee V=1\) and \(U\ne0\). Conversely those two
conditions permit \(Z_{10}=U\), \(Z_{11}=\neg U\); the cover condition
gives \(\neg U\le V\). Set the other cells to zero. This proves
equivalence of the Boolean conditions in any Boolean algebra. The preceding
full reduction proof constructs the corresponding witnesses in each factor.
For an empty index set \(U=0\), so the condition is false, correctly
rejecting the negated all-coordinate predicate in the singleton product.
\end{proof}
\end{annexitem}

\begin{annexitem}{annex:p03-fv-power-structure}{thm:p03-fv-power-structure}{Equality, finite refinement, and all factor-formula values}
\begin{proof}
Reflexivity holds because unequal labels in one presentation have disjoint
coefficients; symmetry is immediate. For transitivity write the three
presentations as \((b_r,a_r)\), \((c_t,d_t)\), and \((e_v,z_v)\), with
\(b_r,c_t,e_v\in B\) and \(a_r,d_t,z_v\in A_s\). For \(a_r\ne z_v\),
\(b_r\wedge e_v=\bigvee_t(b_r\wedge c_t\wedge e_v)\).
For each \(t\), either \(a_r\ne d_t\) or \(d_t\ne z_v\), so one
of the assumed presentation equalities makes that summand zero. This proves
transitivity using only finite Boolean operations.

Intersections of one coefficient from each input partition are pairwise
disjoint and their finite join is \(1\), by distributivity. They are thus
a legitimate common refinement. Refine further with alternative presentations
of those same inputs. On a nonzero resulting cell, presentation equality
forces each old label to equal its corresponding new label. A function
therefore gives identical output labels there, and a relation has identical
truth there. Joining the zero mismatching cells proves operation congruence
and relation invariance. A constant is the single-cell construction. All
sorts are nonempty because any label in the nonempty base sort gives such a
constant presentation. In a trivial Boolean algebra the empty presentation
is also available and all presentations are equivalent.

For a factor formula \(\psi\), refine its finite parameter tuple.
On each cell use actual satisfaction in \(A\), including its ordinary
quantifier semantics. Changing presentations or refining cells changes no
label tuple on any nonzero overlap. Its \(\psi\)-truth is therefore the
same on every such overlap, and distributivity shows that the joined truth
value is unchanged. This proves existence and well-definedness for arbitrary
factor formulas, not merely atomic ones. In particular the truth values of
the finite bit-pattern formulas form a Boolean partition, and original
component values are their indicated subjoins.
\end{proof}
\end{annexitem}

\begin{annexitem}{annex:p03-fv-power-reduction}{thm:p03-fv-power-reduction}{Finite witnesses on each nonzero refined cell}
\begin{proof}
The atomic case follows from the operation/refinement construction and
term induction, just as for coordinate products: a relation or equality
holds exactly when its atomic value is \(1_B\). Logical falsity uses
\(0\ne0\). Negation and conjunction follow from induction on the
matrix as in the full-product proof. It remains to justify the same
existential Boolean clause with witnesses lying in this bounded power.

For an existing power witness, refine its presentation with the parameter
presentations. Give each cell the truth-bit pattern of its labels in \(A\).
Joining the cells with a pattern gives \(Z_\sigma\); their joins form a
partition, each lies under the value of \(\eta_\sigma\), and the
\(\psi_j\)-values are the indicated pattern subjoins. The inductive
matrix condition proves the forward implication.

Conversely suppose \((Z_\sigma)_\sigma\) satisfies the Boolean condition.
Let \((c_t,\bar a_t)_t\) be the common finite parameter partition; for
no parameters use the single cell \(1_B\) labeled by the empty tuple.
Refine by \(Z_\sigma\wedge c_t\). On a nonzero such cell,
\(A\models\eta_\sigma(\bar a_t)\). Indeed, if that sentence were false,
the definition of its factor value would make \(c_t\) disjoint from it,
contradicting \(Z_\sigma\le[\![\eta_\sigma]\!]_B\) and nonzeroness.
For each of the finitely many nonzero cells choose an actual
\(w_{\sigma,t}\in A_s\) satisfying
\(\delta_\sigma(\bar a_t,w_{\sigma,t})\). Finite choice follows in
classical ZF by induction on the size of this finite family.
Omit zero cells and use the remaining coefficients and these labels as the
presentation of \(b\). They still join to \(1_B\). When \(B\) is
trivial all cells are zero and the allowed empty presentation is the witness.
Its \(\psi_j\)-values are precisely \(\bigvee_{\sigma_j=1}Z_\sigma\),
since every nonzero refined cell realizes all bits of its pattern.
The induction hypothesis now proves the matrix at \((\bar a,b)\).
This completes the existential step and hence all formulas. Neither
ultrafilter-indexed Choice nor joins of infinitely many Boolean elements
appear in this proof.
\end{proof}
\end{annexitem}

\begin{annexitem}{annex:p03-fv-power-representations}{thm:p03-fv-power-representations}{Represented fibers and compact Stone spaces}
\begin{proof}
For \(\mathcal B\subseteq\mathcal P(I)\), evaluate a finite labeled
partition at an index using its unique cell. It gives a finite-range
function, and each fiber is a finite union of its coefficient sets in
\(\mathcal B\). Equivalent presentations give the same function.
Conversely, equal functions have no index in a mismatching cell intersection,
so the presentations are equivalent. A finite-range function with fibers
in \(\mathcal B\) is presented by its finitely many actual values and
their fibers. These maps are inverse. The refinement descriptions make them
commute with all operations, and all-index relation truth agrees with the
condition that its joined value is \(1\). Thus they give an isomorphism,
including the unique empty function when \(I\) is empty.

If \(X\) is compact and a function to discrete \(A_s\) is continuous,
its fibers are open, and their complements, unions of the other fibers,
are open too. The fibers are clopen and cover \(X\). Compactness selects
finitely many covering fibers, proving the range finite; no selection of a
neighborhood at every point is needed. Conversely, finite clopen fibers
make the preimage of every subset of the discrete carrier open, proving
continuity. The empty-space case has empty range.

Under BPIT, the already supplied Stone theorem gives an isomorphism
\(B\to\operatorname{Clop}(X_B)\). Transport every finite partition
coefficient along that isomorphism, apply the preceding represented-space
construction, and use the continuity criterion. Coefficient transport
preserves presentation equality, refinements and relation values; this is
the asserted isomorphism to \(A[B]^*\). The relational extension is
exactly the all-point interpretation just checked, not an additional
unproved theorem about arbitrary Boolean products. For trivial \(B\),
Stone space is empty and the construction gives the unique empty function
in every sort, as stated.
\end{proof}
\end{annexitem}

\begin{annexitem}{annex:p03-fv-power-elementarity}{thm:p03-fv-power-elementarity}{Using the computed Boolean conditions}
\begin{proof}
For \(A\preccurlyeq A'\), inclusion of labels preserves and reflects
their equality, so presentation equivalence and every finite refinement
give an injective structure map. For each factor formula \(\psi_j\),
elementarity makes its truth at labels from \(A\) identical in \(A'\).
Thus its Boolean value is unchanged. The power-reduction theorem makes
every product-language formula agree on these parameters, proving the
claimed elementary inclusion.

For \(B\preccurlyeq C\), the Boolean embedding preserves and reflects
zero intersections and preserves all finite joins. It therefore gives an
injective structure map on presentations. Each factor-formula value is
sent to its coefficient image. Elementarity of the Boolean embedding makes
the formula \(\Theta_\varphi\) agree on these images, and the reduction
theorem proves elementarity of the power map.
For a sentence, the component formulas have no free variables. Each such
component value is \(1\) or \(0\) according to its truth in the base
structure. Elementary equivalence of \(A,A'\) makes those choices the
same; elementary equivalence of \(B,B'\) then makes the resulting
Boolean sentence agree. This proves the third assertion, without identifying
the base structures or their Boolean algebras.

Finally, for \(B=\mathcal P(I)\), identify the partition power with
finite-range functions. On such a parameter tuple, its factor truth sets
are literally the same subsets as those in the full power. The direct and
bounded power theorems use the identical Boolean formula in the identical
index algebra, so all first-order formulas agree. This proves the
finite-range substructure elementary in the full power under the full-product
standing, not the unjustified constant-diagonal assertion.
\end{proof}
\end{annexitem}

\begin{annexitem}{annex:p03-fv-patching-characterization}{thm:p03-fv-patching-characterization}{Constants, equalizers, and finite binary mixing}
\begin{proof}
Continuous functions have finite range and clopen fibers by the compactness
argument above. Their equality set is a finite union of intersections of
corresponding fibers, hence clopen. Constants are continuous, and binary
patching along a clopen is continuous because each resulting fiber is
\((f^{-1}(a)\cap C)\cup(g^{-1}(a)\cap(X\setminus C))\).

Conversely let \(S\) satisfy the three conditions. For \(f\in S\)
and \(a\in A\), its fiber is the equality set of \(f\) with the
constant at \(a\), so is clopen. Thus \(f\) is continuous and, by
compactness, finite-range. This proves \(S\) is contained in the bounded
power. For any function of the bounded power, its finitely many nonempty
clopen fibers partition \(X\). Start with a constant and patch the required
constant value on one fiber at a time. Closure under binary patching yields
that function in \(S\). For an empty space, every constant is the unique
empty function, so the first condition already proves equality. This finishes
both inclusions, with no infinite patching step.
\end{proof}
\end{annexitem}
\begin{annexitem}{annex:p03-fv-direct-reduction}{thm:p03-fv-direct-reduction}{Every recursive clause and the actual existential witnesses}
\begin{proof}
Induct on the finite syntax tree of \(\varphi\). The preceding coordinate
lemma proves the atomic clause. Falsity is never satisfied, and \(0\ne0\)
is never satisfied in any Boolean algebra, including the trivial one.
For negation, the induction hypothesis identifies satisfaction of its
subformula with the Boolean condition; logical negation of both sides gives
exactly the specified clause. For conjunction, use disjointly renamed lists
and the two induction hypotheses. The clauses for other connectives and
universal quantification are their declared syntactic expansions.

Suppose first that a product witness \(b\) satisfies \(\chi(\bar a,b)\).
At each index \(i\), the finite list of truth values of
\(\psi_j(\bar a(i),b(i))\) has exactly one bit pattern \(\sigma\).
Let \(Z_\sigma\) be the set of indices with that pattern. These sets form
a finite partition, with zero cells allowed. Membership in \(Z_\sigma\)
supplies the local witness \(b(i)\) for \(\eta_\sigma\), so
\(Z_\sigma\subseteq[\![\eta_\sigma(\bar a)]\!]\).
Also \([\![\psi_j(\bar a,b)]\!]=
\bigcup_{\sigma_j=1}Z_\sigma\). The induction hypothesis for \(\chi\)
therefore gives the last condition of \(\Theta_\exists\). This proves
the forward existential implication.

Conversely let the Boolean formula supply a partition \((Z_\sigma)_\sigma\)
with the specified inclusions and condition. For each index \(i\), its unique
cell \(Z_{\sigma(i)}\) lies inside the local truth set of
\(\exists y_s\delta_{\sigma(i)}\). Hence the set
\[
 W_i:=\{d\in(M_i)_s:
       M_i\models\delta_{\sigma(i)}(\bar a(i),d)\}
\]
is nonempty. The Axiom of Choice chooses \(b(i)\in W_i\) for all indices,
giving an actual element of the quantified product sort. Because the pattern
specifies every positive and negative literal, its truth pattern is exactly
\(\sigma(i)\), not merely a subset of the requested true entries.
Consequently its local \(\psi_j\)-sets are exactly the finite unions in
\(\Theta\). The induction hypothesis now gives
\(M\models\chi(\bar a,b)\), hence the existential formula.

For an empty index set all partition cells are empty and \(b\) is the
unique empty function; the same reasoning applies. When \(m=0\), its one
empty pattern requires only a witness in each nonempty factor sort; the
empty conjunction has no additional truth demands. Each recursion step
replaces a proper subformula by finitely many formulas and Boolean terms.
The bit-pattern family is finite, so induction also proves termination and
effective finite construction. It provides no algorithm for deciding
satisfaction in unknown structures.
\end{proof}
\end{annexitem}
